\documentclass[11pt]{beamer}
\usetheme{Madrid}
\usecolortheme{seahorse}
\usepackage[utf8]{inputenc}
\usepackage[T1]{fontenc}
\usepackage[spanish]{babel}

\title{Quibdó Seguro}
\subtitle{Plataforma de reportes ciudadanos y recompensas}
\author{Equipo de Desarrollo}
\date{2026}

\begin{document}
\frame{\titlepage}

\begin{frame}{Resumen}
\begin{itemize}
  \item Monolito UI (Laravel) + Microservicios
  \item Base de datos: MongoDB
  \item Objetivo: facilitar reportes ciudadanos y gestionar recompensas
\end{itemize}
\end{frame}

\begin{frame}{Microservicios}
\begin{itemize}
  \item Auth Service (tokens, usuarios)
  \item Incident Service (reportes)
  \item Rewards Service (ofertas y canjes)
  \item Notification Service (notificaciones)
  \item Analytics Service (métricas)
  \item API Gateway (8000)
\end{itemize}
\end{frame}

\begin{frame}{Despliegue}
\begin{itemize}
  \item Desarrollo: `docker-compose` con MongoDB y servicios
  \item Producción: Kubernetes / PaaS
\end{itemize}
\end{frame}

\begin{frame}{Accesibilidad y UX}
\begin{itemize}
  \item Colores agradables y legibles
  \item Diseño responsivo (móvil y escritorio)
  \item Manual de usuario para públicos mayores (incluido)
\end{itemize}
\end{frame}

\begin{frame}{Siguientes pasos}
\begin{itemize}
  \item Compilar y generar PDFs
  \item Preparar hosting gratuito y guía paso a paso
  \item Entrega final: PDFs + repositorio funcionando
\end{itemize}
\end{frame}

\end{document}
